\section{Boolean statistics}\label{sec:booleanStat}

%\red{We here study the face CP ranks in case of boolean statistics.
%We further show that any elementary \ComputationActivationNetwork{} to boolean statistics is a maximum entropy distribution.}

We say that a statistic is boolean, if $\headdimof{\selindex}=2$ for all $\selindexin$.
The mean polytope is then a 0/1-polytope \cite{ziegler_lectures_2000}, which vertices are exactly the contained boolean vectors.
The special case of boolean statistics is especially interesting for two reasons.
First, we can provide additional intuition on those faces representable by elementary face activating tensors, namely based on intersections with cube faces.
Based on these properties we also show that any elementary \ComputationActivationNetwork{} to a boolean statistic is a maximum entropy distribution.
Second, the statistic itself can be interpreted in terms of logical formulas.
We exploit this property to provide examples and to characterize each 0/1-polytope as a mean polytope to a collection of propositional formulas.

%For boolean statistics $\hlnstat:\facstates\rightarrow\bigtimes_{\selindexin}[2]$ the mean polytope is a subset of the cube $\fullparcube$.
%% NEEDED?
%In this case, any boolean vector in $\meansetof{\hlnstat,\basemeasure}$ is a vertex.
%It follows, that any distribution reproducing a mean parameter $\meanparamwith$ on the relative interior of $\meansetof{\hlnstat,\basemeasure}$ is positive with respect to $\basemeasure$.

%We apply the exponential distribution characterization of the maximum entropy distribution and get that the maximum entropy distribution is in $\elrealizabledistsof{\sstat}$, if and only if the face measure is in $\elrealizabledistsof{\sstat}$.
%This is exactly the case, when the face is an intersection of the mean polytope with a face of the cupe $\fullparcube$.

\subsection{Characterization of elementary representable faces}

Orienting on \exaref{exa:hypercubeFaces} we define cube-likeness of faces and polytopes.

\begin{definition}
    \label{def:cubeLike}
    We say that a face $\facesymbol$ of $\genmeanset$ is cube-like, if it is empty or there is $\arbset\subset[\seldim]$ and $\headindexof{\arbset}\in\bigtimes_{\selindex\in\arbset}[2]$ such that
    \begin{align*}
        \facesymbol
        = \genmeanset \cap \facesymbolof{(\arbset,\headindexof{\arbset})} \, .
    \end{align*}
    Here we denote by $\facesymbolof{(\arbset,\headindexof{\arbset})}$ the faces of the hypercube $\cubeof{\seldim}$ (see \exaref{exa:hypercubeFaces}).
    We further say that a polytope $\genmeanset$ is cube-like, if all faces $\facesymbol\in\facelatticeof{\genmeanset}$ are cube-like.
\end{definition}

We now show that a face is cube-like if and only if it is representable by an elementary tensor.

\begin{theorem}
    \label{the:faceMeasureHardLogicNetworks}
    Let $\hlnstat$ be a boolean statistic, $\basemeasure$ a base measure and $\facesymbol$ be a face of $\genmeanset$.
    Then the following are equivalent:
    \begin{itemize}
        \item[(i)] $\facesymbol$ is cube-like (see \defref{def:cubeLike}).
        \item[(ii)] $\facesymbol$ is representable by an elementary tensor (see \defref{def:faceRepresentability}).
    \end{itemize}
\end{theorem}
\begin{proof}
    If the face is empty, i.e. $\facesymbol=\varnothing$, it is by definition cube-like and has $\zerosat{\headvariables}$ as an elementary activation tensor.
    We therefore assume in the following $\facesymbol\neq\varnothing$.

    (i)$\Rightarrow$(ii):
    Let us assume that $\facesymbol$ is cube-like, that is there is $\arbset\subset[\seldim]$ and $\headindexof{\arbset}\in\bigtimes_{\selindex\in\arbset}[2]$ such that $\facesymbol = \genmeanset \cap \facesymbolof{(\arbset,\headindexof{\arbset})}$.
    We use that $\facesymbol,\genmeanset$ and $\facesymbolof{(\arbset,\headindexof{\arbset})}$ are the convex hulls of the cube vertex sets $\imsetof{\hlnstat,\meanset}{\facesymbol}$, $\imsetof{\hlnstat,\meanset}{}$ and
    \begin{align}\label{eq:cubeFaceVertices}
        \imsetof{}{\arbset,\headindexof{\arbset}}
        = \{\imelementwith \wcols \uniquantwrtof{\selindex\in\arbset}{\imelementat{\indexedselvariable}=\headindexof{\arbset}}\} \, ,
    \end{align}
    which implies that
    \begin{align*}
        \imsetof{\hlnstat,\meanset}{\facesymbol}
        = \imsetof{}{\arbset,\headindexof{\arbset}} \cap \imsetof{\hlnstat,\meanset}{} \, .
    \end{align*}
    Thus, we can choose $\arbset=\imsetof{\hlnstat,\meanset}{\facesymbol}$ for the representation of the face $\facesymbol$ (see \defref{def:faceRepresentability}) and further have that
    \begin{align*}
        \sum_{\imelementwith\in\imsetof{}{(\arbset,\headindexof{\arbset})}}\onehotmapofat{\imelementwith}{\headvariables}
        = \bigotimes_{\selindexin} \hardactlegwith \, .
    \end{align*}
    We have thus found an elementary activation tensor for the face $\facesymbol$.

    (ii)$\rightarrow$(i)
    Conversely, let $\acttensorwith$ be an elementary activation tensor of the face $\facesymbol$ in $\genmeanset$.
    Since $\acttensorwith$ is the sum of different one-hot encodings it is boolean and we find an elementary decomposition $\acttensorwith=\bigotimes_{\selindexin}\acttensorlegwith$ such that the leg vectors $\acttensorwith$ are boolean.
    Since $\facesymbol\neq\varnothing$ we further have for $\selindexin$ that $\acttensorlegwith\neq\zerosat{\headvariableof{\selindex}}$, and thus $\acttensorlegwith\in \{\fbasisat{\headvariableof{\selindex}},\tbasisat{\headvariableof{\selindex}},\onesat{\headvariableof{\selindex}}\}$
    We construct a set $\arbset\subset[\seldim]$
    \begin{align*}
        \arbset
        = \left\{\selindex \wcols {\acttensorlegwith}\neq\onesat{\headvariableof{\selindex}}\right\}
    \end{align*}
    and a tuple
    \begin{align*}
        \headindexof{\selindex}
        = \begin{cases}
              1 & \ifspace {\acttensorlegwith}=\onehotmapofat{1}{\headvariableof{\selindex}} \\
              0 & \ifspace {\acttensorlegwith}=\onehotmapofat{0}{\headvariableof{\selindex}}
        \end{cases} \, .
    \end{align*}
    By construction ${\acttensorlegwith} = \hardactlegwith$ (see \exaref{exa:hypercubeFaces}) follows for all $\selindexin$ and therefore $\acttensorwith=\hardacttensorwith$.
    We therefore have for the set \eqref{eq:cubeFaceVertices}
    \begin{align*}
        \acttensorwith
        = \sum_{\imelement\in\imsetof{}{\arbset,\headindexof{\arbset}}} \onehotmapofat{\imelement}{\headvariables}
    \end{align*}
    and since $\acttensorwith$ is an activation tensor for $\facesymbol$, that
    \begin{align*}
        \imsetof{\hlnstat,\meanset}{\facesymbol}
        = \imsetof{}{\arbset,\headindexof{\arbset}}\cap\imsetof{\hlnstat,\meanset}{} \, .
    \end{align*}
    Taking convex hulls on both sides, this is equivalent to
    \begin{align*}
        \facesymbol
        = \facesymbolof{(\arbset,\headindexof{\arbset})} \cap \genmeanset \, .
    \end{align*}
    We conclude that $\facesymbol$ is cube-like.
\end{proof}

%% Reference to simplices and hypercubes
\theref{the:faceMeasureHardLogicNetworks} implies that the maximum entropy distributions to hypercubes (see \exaref{exa:hypercubeFaces}) and simplices (see \exaref{exa:simplexFaces}) are in $\realizabledistsof{\hlnstat,\elgraph,\basemeasure}$, since these polytopes are cube-like.

% Theorem intro
We can further show, that for boolean statistics, any distribution in $\realizabledistsof{\hlnstat,\elgraph,\basemeasure}$ is a maximum entropy distribution.

%\subsection{Set of maximum entropy distributions}

\begin{theorem}
    Let $\hlnstat$ be a boolean statistic.
    Any distribution in $\realizabledistsof{\hlnstat,\elgraph,\basemeasure}$ is a maximum entropy distribution with respect to $(\hlnstat,\meanparam,\basemeasure)$ where $\meanparam$ is its mean parameter.
    %Any maximum entropy distribution is realized by $\realizabledistsof{\hlnstat,\elgraph,\basemeasure}$ if and only if the mean parameter is in the relative interior of a cube-like face.
\end{theorem}
\begin{proof}
    %To show the first claim,
    Let $\probtensor\in\realizabledistsof{\hlnstat,\elgraph,\basemeasure}$ be arbitrary.
    We then find an elementary tensor $\bigotimes_{\selenumeratorin}\hypercoreofat{\selindex}{\headvariableof{\selindex}}$ and $\partitionfunction\in\rr$ such that
    \begin{align*}
        \probwith
        = \frac{1}{\partitionfunction} \contractionof{\bencsstatwith,\bigotimes_{\selenumeratorin}\hypercoreofat{\selindex}{\headvariableof{\selindex}}}{\shortcatvariables} \, .
    \end{align*}
    We then define for each $\selindexin$ the vector $\kcoreofat{\selindex}{\headvariableof{\selindex}}$ as the support of $\hypercoreofat{\selindex}{\headvariableof{\selindex}}$ and
    \begin{align*}
        \weightparamat{\indexedselvariable}
        = \begin{cases}
              \lnof{\frac{\hypercoreofat{\selindex}{\headvariableof{\selindex}=1}}{\hypercoreofat{\selindex}{\headvariableof{\selindex}=0}}} & \ifspace \hypercoreofat{\selindex}{\headvariableof{\selindex}=0},\hypercoreofat{\selindex}{\headvariableof{\selindex}=1} \neq 0 \\
              0 \elsetext
        \end{cases} \, .
    \end{align*}
    Since $\hlnstat$ is a boolean statistic the $\bigotimes_{\selindexin}\kcoreofat{\selindex}{\headvariableof{\selindex}}$ is a face activation tensor to the face
    \begin{align*}
        \facesymbol \coloneqq \hlnmeanset \cap \facesymbolof{(\arbset,\headindexof{\arbset})}
    \end{align*}
    where $\arbset=\{\selindex\wcols\kcoreofat{\selindex}{\headvariableof{\selindex}}\neq \onesat{\headvariableof{\selindex}}\}$ and for $\selindex\in\arbset$
    \begin{align*}
        \headindexof{\selindex}
        = \begin{cases}
              0 & \ifspace \kcoreofat{\selindex}{\headvariableof{\selindex}=1} = 0 \\
              1 & \elsetext
        \end{cases} \, .
    \end{align*}
    Therefore, $\probtensor$ is the maximum entropy distribution to the canonical parameters $\facesymbol$ and $\weightparam$.
%    First claim by decomposing any elementary tensor into exponential and hard activation core.
%    Second claim by characterization of elementary faces by cube-likeness.
\end{proof}

\subsection{Interpretation by propositional formulas}

We can understand each feature as a propositional formula and the variables $\shortcatvariables$ as atoms (possibly after a binarization).
We use propositional connectives for syntactical representations of the formulas and denote by $\lnot^{1}$ the logical negation and by $\lnot^{0}$ the identity.
A boolean tensor $\formulaat{\shortcatvariables}$ is called satisfiable, if there is at least one $\shortcatindicesin$ with $\formulaat{\indexedshortcatvariables}=1$.

%% Vertex characterization
%Each vertex of the cube, which is not a vertex $\imelementwith$ of the polytope corresponds with the unsatisfiability of a formula
%\begin{align*}
%    \bigwedge_{\selindexin} \lnot^{1-\imelementat{\indexedselvariable}} \formulaofat{\selindex}{\shortcatvariables}
%\end{align*}
%which is equal with any of the entailment statements for $\arbset\subset[\seldim]$ % Can extend to any partition!
%\begin{align*}
%    \left(\bigwedge_{\selindex\in\arbset} \lnot^{1-\imelementat{\indexedselvariable}} \formulaofat{\selindex}{\shortcatvariables} \right)
%    \models \left(\bigwedge_{\selindex\in\arbset}\lnot^{\imelementat{\indexedselvariable}} \formulaofat{\selindex}{\shortcatvariables} \right) \, .
%\end{align*}

\begin{theorem}\label{the:boolVecSat}
    Let $\sstat$ be a boolean statistic.
    Any boolean vector $\imelementwith\in\{0,1\}^{\seldim}$ is a vertex of $\hlnmeanset$ if and only if
    \begin{align*}
        \formulaofat{\imelement}{\shortcatvariables}
        \coloneqq \bigwedge_{\selindexin} \lnot^{1-\imelementat{\indexedselvariable}} \formulaofat{\selindex}{\shortcatvariables}
    \end{align*}
    is satisfiable.
\end{theorem}
\begin{proof}
    Let $\imelementwith\in\{0,1\}^{\seldim}$ be arbitrary.
    If and only if $\imelementwith$ is a vertex of $\hlnmeanset$ we find $\shortcatindicesin$ with $\imelementwith=\sencsstat{\indexedshortcatvariables,\selvariable}$.
    This is further equivalent to $\formulaofat{\selindex}{\indexedshortcatvariables}=\imelementat{\indexedselvariable}$ for any $\selindexin$, $\lnot^{1-\imelementat{\indexedselvariable}}\formulaofat{\selindex}{\indexedshortcatvariables}=1$ for any $\selindexin$, and $\formulaofat{\imelement}{\indexedshortcatvariables}=1$.
\end{proof}

%% Entailment interpretation
Since the presence of boolean vectors in $\hlnmeanset$ is by \theref{the:boolVecSat} related to the satisfiability of the conjunction $ \formulaofat{\imelement}{\shortcatvariables}$, the absence is related to entailment statements.
To be more precise, for any $\imelementwith\in\{0,1\}^{\seldim}$ and arbitrary $\arbset\subset[\seldim]$ we have $\imelementwith\notin\hlnmeanset$, if and only if
\begin{align*}
    \left(\bigwedge_{\selindex\in\arbset} \lnot^{1-\imelementat{\indexedselvariable}} \formulaofat{\selindex}{\shortcatvariables} \right)
    \models \left(\bigvee_{\selindex\notin\arbset}\lnot^{\imelementat{\indexedselvariable}} \formulaofat{\selindex}{\shortcatvariables} \right) \, .
\end{align*}
By definition (see \cite{russell_artificial_2021}) the latter holds if and only if
\begin{align*}
    \left(\bigwedge_{\selindex\in\arbset} \lnot^{1-\imelementat{\indexedselvariable}} \formulaofat{\selindex}{\shortcatvariables} \right)
    \land \lnot \left(\bigvee_{\selindex\notin\arbset}\lnot^{\imelementat{\indexedselvariable}} \formulaofat{\selindex}{\shortcatvariables} \right)
    = \formulaofat{\imelement}{\shortcatvariables}
\end{align*}
is not satisfiable.

Atomic formulas and minterm formulas correspond to hypercubes and simplices as we show next.

\begin{example}[Atomic formulas]
    \label{exa:atomicFormulasHypercube}
    %The assumption of \theref{the:sufficientHLNExpressivity} is satisfied in
    Let us consider the case of atomic formulas $\formulaofat{\selindex}{\shortcatvariables}$ defined with the coordiantes at $\shortcatindicesin$ by
    \begin{align*}
        \formulaofat{\selindex}{\indexedshortcatvariables}
        = \catindexof{\selindex} \, .
    \end{align*}
    The mean polytope in this case is the $\catorder$-dimensional hypercube (see \exaref{exa:hypercubeFaces})
    \begin{align*}
        \meansetof{\atomformulaset,\ones}
        = \fullparcube
    \end{align*}
    which is called a simple polytope, since each vertex is contained in the minimal number of $\catorder$ facets.
    %Since the cube $\fullparcube$ is a face of itself, \theref{the:sufficientHLNExpressivity} implies $\hlnmeanset|_{\hlnsetof{\atomformulaset}}=\hlnmeanset$.

    Each face is characterized by the projections onto each variable, which is either $\{0\}$, $\{1\}$ or $[0,1]$.
    The projections are represented by the tuple $\hardparam$ defined in the following way:
    \begin{itemize}
        \item We define the set $\hardlegset\subset[\atomorder]$ of variables, such that the projection onto the variable is $\{0\}$ or $\{1\}$
        \item We define to each $\selindex\in\hardlegset$ an index $\headindexof{\selindex}=0$ if the projection is $\{0\}$ and $\headindexof{\selindex}=1$ if the projection is $\{1\}$.
    \end{itemize}
    Trivially, each face of the hypercube is a cube face. %$\elmeansetof{\atomformulaset}=\meansetof{\atomformulaset}$.

    % Extension: Discussion of those statistics having the hypercube as mean polytope
    More general, the mean polytope to a collection $\formulaofat{\selindex}{\shortcatvariables}$ of propositional formulas is the hypercube, if and only if for any boolean vector $\imelementwith$ we have
    \begin{align*}
        \contraction{\bigwedge_{\selindexin}\lnot^{1-\imelementat{\indexedselvariable}}\formulaofat{\selindex}{\shortcatvariables}}
        \neq 0 \, .
    \end{align*}
    This means that no conjunction of possibly negated formulas from this statistic entails nor contradicts another formula.

\end{example}

\begin{example}[Minterm formulas build simplices]
    \label{exa:mintermHLNSet}
    The set of minterm formulas is indexed by $\imelementwith\in\{0,1\}^{\seldim}$ and given by
    \begin{align*}
        \formulaofat{\imelement}{\shortcatvariables}
        = \bigwedge_{\selindex\in[\catorder]} \lnot^{1-\imelementat{\indexedshortcatvariables}} \, \catvariableof{\selindex}
        = \onehotmapofat{\imelementwith}{\shortcatvariables}
    \end{align*}
    where by $\catvariableof{\selindex}$ we denote the $\selindex$-th atomic formula (see \exaref{exa:atomicFormulasHypercube}).
    The mean polytope to the minterm statistic and base measure $\basemeasure$ is the standard simplex of dimension
    \begin{align*}
        \cardof{\left\{\imelementwith\wcols\imelement\in\{0,1\}^{\seldim}\ncond\basemeasureat{\shortcatvariables=\imelementwith}\neq 0\right\}}-1 \, ,
    \end{align*}
    that is the set
    \begin{align*}
        \meansetof{\mintermformulaset,\basemeasure}
        = \convhullof{\onehotmapofat{\imelementwith}{\headvariables}\wcols\imelement\in\{0,1\}^{\seldim}\ncond\basemeasureat{\shortcatvariables=\imelementwith}\neq 0} \, .
    \end{align*}
    We notice that for boolean base measures $\meansetof{\mintermformulaset,\basemeasure}=\realizabledistsof{\basemeasure}$.
    % Cube like
    The simplex is further cube-like, since each face is indexed by a subset $\arbset\subset\{0,1\}^{\seldim}$ of vertices and the intersection of the simplex with the cube face
    $\facesymbolof{(\{0,1\}^p\backslash\arbset,\onetuple{\{0,1\}^p\backslash\arbset})}$.

    %of is the intersection of $\hlnsetof{\mintermformulaset}$ with the cube face $\facesymbolof{\hardparam}$.

%    In this case, $\hlnsetof{\mintermformulaset}$ contains any distribution and therefore trivially realizes any mean parameter in $\meansetof{\mintermformulaset,\ones}$.

%    The faces of the standard simplex are itself standard simplices to base measures $\secbasemeasure$ with $\secbasemeasure\prec\basemeasure$.
%    We store them by the tuple $\hardparam$, where $\hardlegset$ is the support of $\secbasemeasure$ and $\headindexof{\selindex}=0$ for $\selindex\in\hardlegset$.


    % 
    More general the mean polytope to a boolean statistic is a standard simplex, if and only if each target coordinate restriction $\formulaof{\secselindex}$ contradicts all others, that is
    \begin{align*}
        \uniquantwrtof{\selindex\neq\secselindex\in[\seldim]}{\formulaof{\selindex}\models\lnot\formulaof{\secselindex}} \, .
    \end{align*}
\end{example}

Let us now provide examples of statistics, which polytopes are not cube-like.

\input{examples/boolean-stat/nonel_hlnstat_maxent}

%% Order?
\input{examples/boolean-stat/tt_diagonal_face}

\subsection{Construction of boolean statistics to polytopes}

We now show, that any convex polytope with boolean vertices in $\parspace$ is the mean polytope to a boolean statistic.
To construct the corresponding formulas we denote by $\restrictionofto{\meanset}{\arbset}$ the restriction of the vectors in the set $\meanset$ to the coordinates in $\arbset$.

\begin{theorem}\label{the:boolStatToPolytope}
    Let $\meanset$ an arbitrary polytope with boolean vertices in $\parspace$.
    Then we construct propositional formulas on atoms $\catvariableof{[\seldim]}$ by
    \begin{align*}
        \formulaofat{0}{\catvariableof{0}} =
        \begin{cases}
            \top & \ifspace \restrictionofto{\meanset}{\{0\}} = \{1\} \\ %\restrictionofto{\meanset}{\rr^1\times 0_{\seldim-1}} = \{1\} \\
            \bot & \ifspace \restrictionofto{\meanset}{\{0\}} = \{0\} \\%\restrictionofto{\meanset}{\rr^1 \times 0_{\seldim-1}} = \{0\} \\
            \catvariableof{0} & \ifspace \restrictionofto{\meanset}{\{0\}} = [0,1] %\restrictionofto{\meanset}{\rr^1 \times 0_{\seldim-1}} = [0,1]
        \end{cases}
    \end{align*}
    and iteratively for $\selindexin$ with $\selindex\geq1$ by
    \begin{align*}
        \formulaofat{\selindex}{\catvariableof{[\selindex+1]}} =
        \bigwedge_{\imelementwith \in \restrictionofto{(\meanset\cap \{0,1\}^{\seldim})}{[\selindex]}}
        \left(\left(\bigwedge_{\secselindex\in[\selindex]} \lnot^{1-\imelementat{\selvariable=\secselindex}} \formulaofat{\secselindex}{\catvariableof{[\secselindex+1]}}\right) \Rightarrow
            \begin{cases}
                \top & \ifspace \restrictionofto{(\meanset\cap\imelement\times\rr^{\seldim-\selindex})}{\{\selindex\}} = \{1\} \\%\restrictionofto{\meanset}{v\times\rr^{1} \times 0_{\seldim-\selindex-1}} = \{1\} \\
                \bot & \ifspace \restrictionofto{(\meanset\cap\imelement\times\rr^{\seldim-\selindex})}{\{\selindex\}} = \{0\} \\ %\restrictionofto{\meanset}{v\times\rr^{1} \times 0_{\seldim-\selindex-1}} = \{0\} \\
                \catvariableof{\selindex} & \ifspace \restrictionofto{(\meanset\cap\imelement\times\rr^{\seldim-\selindex})}{\{\selindex\}} = [0,1] %\restrictionofto{\meanset}{v\times\rr^{1} \times 0_{\seldim-\selindex-1}} = \{0,1\} \\
            \end{cases}
        \right).
    \end{align*}
    Here we denote by $\restrictionofto{\meanset}{V}$ the projections of the vertices in $\meanset$ onto the subspaces $V$, and by $0_{\seldim}$ the zero vector in $\rr^{\seldim}$.
\end{theorem}
\begin{proof}
    We show per induction, that for any $\selindexin$ the family of \HybridLogicNetworks{} with the statistic $\formulaof{[\selindex+1]}$ by the first $\selindex+1$ formulas has the mean polytope
    \begin{align}
        \label{eq:HLNindConTBS}
        \meansetof{\formulaof{[\selindex+1]},\trivbm}
        = \restrictionofto{\meanset}{[\selindex]} \, .
    \end{align}

    $"\selindex=0"$: The polytope $\restrictionofto{\meanset}{[\selindex]}=\{1\}$ (respectively $\restrictionofto{\meanset}{[\selindex]}=\{0\}$) is reproduced by $\formulaof{0}$ being a tautology (respectively a contradiction).
    In the case $\restrictionofto{\meanset}{[\selindex]}=[0,1]$ the polytope is reproduced by the any formula, which is neither a tautology nor a contradiction, and the atomic formula $\catvariableof{0}$ is an example of such an contingency.
    Since the projection of a 0-1 polytope onto the first coordinates is itself a 0-1 polytope, these are the only possible cases and we conclude that in all
    \begin{align*}
        \meansetof{\formulaof{[\selindex]},\trivbm}
        = \restrictionofto{\meanset}{[\selindex]} \, .
    \end{align*}

    $"\selindex\rightarrow\selindex+1"$:
    We define for any $\imelement\in\meansetof{\formulaof{[\selindex]},\trivbm}\cap\{0,1\}^{\selindex}$ the formula
    \begin{align*}
        \formulaofat{\selindex,\imelement}{\catvariableof{[\selindex+1]}}\coloneqq
        \bigwedge_{\secselindex\in[\selindex]} \lnot^{1-\imelementat{\selvariable=\secselindex}} \formulaofat{\secselindex}{\catvariableof{[\secselindex+1]}} \, .
    \end{align*}
    and notice that all formulas are satisfiable (since the corresponding vertices are in the $\meansetof{\formulaof{[\selindex]},\trivbm}$) and have disjoint models.
    We have for any boolean vector $\imelement\in\{0,1\}^{\selindex}$ and $\lambda\in\{0,1\}$ that $\imelement\times\{\lambda\}\in\restrictionofto{\meanset}{[\selindex+1]}$, if and only if
    \begin{align*}
        \imelement \in \restrictionofto{\meanset}{[\selindex]} \quad \text{and} \quad
        \lambda \in \restrictionofto{\meanset\cup\imelement\times\rr^{\seldim-\selindex}}{[\selindex]} \, .
    \end{align*}
    With modus ponens on $\formulaofat{\selindex+1}{\catvariableof{[\selindex+2]}}$ this is further equivalent to the satisfaction of
    \begin{align*}
        \formulaofat{\selindex,\imelement}{\catvariableof{[\selindex+1]}} \land \lnot^{1-\lambda}\formulaofat{\selindex+1}{\catvariableof{[\selindex+2]}}
    \end{align*}
    and thus to $\imelement\times\{\lambda\}\in\meansetof{\formulaof{[\selindex+1]},\trivbm}$.
    The vertices of the 0/1 polytopes $\meansetof{\formulaof{[\selindex+1]},\trivbm}$ and $\restrictionofto{\meanset}{[\selindex+1]}$ thus coincide and both polytopes are therfore equivalent.

%    $\selindex\rightarrow\selindex+1$: Let us assume, that \eqref{eq:HLNindConTBS} holds for a $\selindexin$.
%    Then for any $\shortcatindicesin$ there is exactly one $\imelementwith \in \restrictionofto{\meanset}{[\selindex]}$ such that $\shortcatindices$ is a model of
%    \begin{align*}
%        \formulaofat{\selindex,\imelement}{\catvariableof{[\selindex+1]}}\coloneqq
%        \bigwedge_{\secselindex\in[\selindex]} \lnot^{1-\imelementat{\selvariable=\secselindex}} \formulaofat{\secselindex}{\catvariableof{[\secselindex+1]}} \, .
%    \end{align*}
%    This holds, since by the mutual contradiction of these formulas at most one can have $\shortcatindices$ as a model.
%    Further, if none would have $\shortcatindices$ as a model, then the to $\shortcatindices$ corresponding vector of satisfactions $(\formulaofat{\selindex}{\indexedshortcatvariables})_{\selindex\in[\selindex]}$ is not in $\restrictionofto{\meanset}{[\selindex]}$, which can not be the case.
%    Now, the implications to all $\tilde{\imelement}$ except for $\imelement$ are for the models $\shortcatindices$ of $\formulaof{\selindex,\imelement}$ satisfied, and the satisfaction of $\formulaof{\selindex}$ thus only depends on the head of the implication to $v$.
%    It follows, that the vertices of $\meansetof{\formulaof{[\selindex+2]},\trivbm}$ sharing the first $\selindex$ coordinates with $\imelement$ are determined by the head of the implication at $\imelement$.
%    With the same arguments as in the case $\selindex=0$ we now notice, that in the three cases we construct vertex sets $\imelement\times\{1\},\, \imelement\times\{0\}$ or $\imelement\times\{0,1\}$ if and only if they appear in the polytope $\restrictionofto{\meanset}{[\selindex+1]}$.
%    This establishes for each vertex $\imelement$ of $\meansetof{\formulaof{[\selindex+1]},\trivbm}$ that
%    \begin{align*}
%        \restrictionofto{(\meansetof{\formulaof{[\selindex+2]},\trivbm}\cap\imelement\times\rr^{\seldim-\selindex})}{[\selindex+1]}
%        = \restrictionofto{(\meanset\cap\imelement\times\rr^{\seldim-\selindex})}{[\selindex+1]} \, .
%    \end{align*}
%    Since for any vertex in $\meanset$ we find a unique vertex in $\meansetof{\formulaof{[\selindex+1]},\trivbm}$ sharing the first $\selindex$ coordinates, we have that \eqref{eq:HLNindConTBS} holds for $\selindex+1$.

    By induction the equation \eqref{eq:HLNindConTBS} holds for arbitrary $\selindexin$.
    For $\selindex=\seldim-1$ this is the claim.
\end{proof}

\subsection{Examples in statistical physics and coding theory}

\input{examples/boolean-stat/ising_two}

\input{examples/boolean-stat/ldpc_one}